\section{Constitutive normalization and spectral incidences of return}
\label{iii:sec:incidencias-espectrales}

The two-sheet response and the return functional preserve two
complementary kinds of information. The first comprises the product
and ratio of the constitutive eigenvalues; the second comprises the
heights, periods and oriented multiplicities of the incidences.
This section assembles the operations that distinguish these readings.
Its antecedent is the angular transport produced by APP--TRIT--TPK,
with its projectors and lifted chart. Normalizations are applied after
constructing that transport and preserving its orientation.

\subsection{Exact separation of the common and oriented modes}

Let $r_+,r_->0$ be the labelled eigenvalues of
\eqref{iii:eq:rchannels}. Define
\begin{equation}
 a=\sqrt{r_+r_-},\qquad
 \eta=\frac12\log\frac{r_+}{r_-},\qquad
 \widetilde{\mathcal V}=a^{-1}\mathcal V.
 \label{iii:eq:rec27-normalizacion}
\end{equation}
In the chamber $x>y>0$, $a\in(3/4,1)$ and $\eta>0$.
The orientation marker $\mathcal R$ satisfies
$\mathcal R P_\pm=\pm P_\pm$.

\begin{proposition}[Positive factorization with preserved orientation]
\label{iii:prop:rec27-factorizacion}
The response admits the unique factorization
\begin{equation}
 \mathcal V=a\exp(\eta\mathcal R),\qquad
 \det\widetilde{\mathcal V}=1.
 \label{iii:eq:rec27-unimodular}
\end{equation}
Its four constitutive readings satisfy
\begin{equation}
 \widehat\varepsilon=a^2e^{2\eta},\quad
 \widehat\mu=a^2e^{-2\eta},\quad
 \widehat Z=e^{-2\eta},\quad
 \widehat c=a^{-2}.
 \label{iii:eq:rec27-cuatro}
\end{equation}
With the orientation fixed, sheet interchange preserves $a$ and sends
$\eta$ to $-\eta$.
\end{proposition}
\begin{proof}
Spectral calculus gives
$\exp(\eta\mathcal R)=e^\eta P_++e^{-\eta}P_-$.
The definitions imply $ae^\eta=r_+$ and $ae^{-\eta}=r_-$,
thus recovering $\mathcal V$ exactly.
The product of its eigenvalues fixes the unique positive root $a$;
their ratio fixes the unique real value of $\eta$.
Substitution into \eqref{iii:eq:four} proves the four identities.
Finally, interchanging the eigenvalues inverts their ratio and preserves
their product.
\end{proof}

Unimodular normalization preserves the entire oriented coordinate and
separates the common coordinate. If $a$ is also retained, the transformation
is reversible. Publishing only $\widetilde{\mathcal V}$ determines
$\widehat Z$, whereas $\widehat c$ requires the archived common factor.
This distinction explains the role of reciprocal sheets in a unit-product
reader: such a reader realizes the unimodular component. The complete
constitutive response additionally preserves its determinant.

In particular, historical readers of the form
$\mu_r^\pm=e^{\pm2u}$ and $\varepsilon_r^\pm=e^{\mp2u}$
have product one on each sheet. Their impedance ratio is $e^{4u}$.
The current reading in \eqref{iii:eq:rec27-cuatro} has product
$a^4=\widehat c^{-2}$. The historical parameter $u$ and the coordinate
$\eta$ are identified only when the map between these realizations is
specified; the preceding decomposition states exactly which additional
factor restores the complete response.

\subsection{Operational role of the Catalan moment}

The reconstruction in
Theorem~\ref{iii:thm:restitucion-funcional-catalan} places Catalan's
constant at the endpoint of a functional collection. The operation entering
the interior response is
\begin{equation}
 r_\pm=
 \frac{1+s_\pm+s_\pm^2}{s_\pm(1+s_\pm)}
 \mathcal D^2\mathcal C_2(s_\pm),\qquad
 s_\pm=q_\pm^{30},\quad \mathcal D=s\frac{d}{ds}.
 \label{iii:eq:rec27-catalan-operativo}
\end{equation}
The collection and its condition at the origin determine the moment;
cubic inversion determines its two arguments from the responses.
Two levels of reconstruction are thus retained: a function from a function,
and arguments from evaluations of an already fixed function.

Composition with \eqref{iii:eq:rec27-normalizacion} obtains $a$ and
$\eta$ from these same evaluations. Subsequent composition with
\eqref{iii:eq:barbero-factorization} obtains the oriented
Barbero--Immirzi functional. The endpoint value $\mathcal C_2(1)$,
the responses $\mathcal D^2\mathcal C_2(s_\pm)$, and the functional
$\mathcal H(s_-)-\mathcal H(s_+)$ therefore perform specific
mathematical roles within the same chain. Reconstruction uses the
complete collection and the uniqueness conditions already proved.

To specify which observation distinguishes the channels, write
$u_\pm=\log r_\pm$. Variations within a common chart satisfy
\begin{equation}
 \begin{pmatrix}d\log\widehat c\\d\log\widehat Z\end{pmatrix}
 =\begin{pmatrix}-1&-1\\-1&1\end{pmatrix}
 \begin{pmatrix}du_+\\du_-\end{pmatrix}.
 \label{iii:eq:rec27-variaciones}
\end{equation}
The determinant is $-2$, so the joint reading recovers both variations.
Each reading separately retains one linear combination.
This identity expresses the complementary roles of speed and impedance,
with orientation and dimensional chart fixed.

\subsection{Three distinct readings of the $90/120$ incidences}

The flags in \eqref{iii:eq:barbero-flags} determine the heights
$90$ and $120$. The normalized defect of their cardinalities is $-1/12$.
The chain in \eqref{iii:eq:barbero-cycle}, in turn, determines the signed
weights $12$ and $-1$. Each height has its return series
$S_m(q)=q^m/(1-q^{3m})$. Preserving these three operations allows an
exact comparison of the two functionals
\begin{equation}
 L_{\rm eq}(q)=S_{90}(q)-S_{120}(q),\qquad
 L_{\rm ret}(q)=12S_{90}(q)-S_{120}(q).
 \label{iii:eq:rec27-dos-kernels}
\end{equation}
The first corresponds to the equal-weight antecedent; the second reads
the current oriented chain.

\begin{proposition}[Normalizations of the return pole]
As $q\uparrow1$,
\begin{equation}
 \lim(1-q)L_{\rm eq}(q)=\frac1{1080},\qquad
 \lim(1-q)L_{\rm ret}(q)=\frac1{24},\qquad
 L_{\rm ret}-L_{\rm eq}=11S_{90}.
 \label{iii:eq:rec27-polos}
\end{equation}
\end{proposition}
\begin{proof}
Factoring $1-q^{3m}$ proves
$\lim(1-q)S_m(q)=1/(3m)$.
By linearity, the two coefficients are respectively
$1/270-1/360=1/1080$ and $12/270-1/360=1/24$.
Subtracting the definitions proves the last identity.
\end{proof}

The ratio $45$ between these singular coefficients follows from the
specified weights. The functions also retain their regular terms;
the ratio of their poles does not identify their values throughout the
interval. The cardinality defect, chain weight and singular coefficient
are thereby assigned distinct operations and normalizations.

\clearpage
\subsection{Fourier resolution of the dodecaphasic pattern}

On the same cycle, the pattern associated with the oriented incidences is
represented by
\begin{equation}
 p_m=12\cos\frac{2\pi\,3m}{12}
      -\cos\frac{2\pi\,4m}{12},\qquad 0\le m<12.
 \label{iii:eq:rec27-patron}
\end{equation}
Frequencies $3$ and $4$ realize phase increments of $90^\circ$ and
$120^\circ$. The arguments of the trigonometric functions are expressed
in radians. The weights come from the previously constructed chain;
the harmonic representation is applied afterwards.

\begin{proposition}[Spectral support and normalization]
\label{iii:prop:rec27-fourier}
For the unnormalized transform
$\widehat p_k=\sum_{m=0}^{11}p_m e^{-2\pi i km/12}$,
\begin{equation}
 \widehat p_3=\widehat p_9=72,\qquad
 \widehat p_4=\widehat p_8=-6,
 \qquad \widehat p_k=0\ \text{at all other indices}.
 \label{iii:eq:rec27-fourier}
\end{equation}
The pattern has zero mean and $\sum_m|p_m|^2=870$.
\end{proposition}
\begin{proof}
Writing each cosine as the half-sum of two exponentials gives coefficients
$6$ at characters $3,9$ and $-1/2$ at $4,8$.
Finite orthogonality,
$\sum_{m=0}^{11}e^{2\pi i(j-k)m/12}=12\delta_{jk}$,
gives \eqref{iii:eq:rec27-fourier}. The coefficient at index zero
vanishes. The same orthogonality gives
$\sum|p_m|^2=(2\cdot72^2+2\cdot6^2)/12=870$.
\end{proof}

For the transform divided by $12$, the coefficients are $6$ and
$-1/2$; under unitary normalization, they are $12\sqrt3$ and $-\sqrt3$.
Their magnitude ratio $12:1$ is preserved in both cases. The exact sequence
\begin{equation}
 2(p_0,\ldots,p_{11})
 =(22,1,-23,-2,25,1,-26,1,25,-2,-23,1)
 \label{iii:eq:rec27-patron-racional}
\end{equation}
allows the transform to be checked with exact algebraic arithmetic.

The zero mean of the pattern comes from character cancellation on the
cycle. The defect $-1/12$ comes from the normalized incidence count;
$1/24$ comes from the pole of the return series. Keeping these domains
specified explains what each reading preserves before composition with
the angular pair and the Barbero--Immirzi functional.
